\subsection{Implementation and System Architecture}

This section describes the software implementation used to realize the prediction methodology introduced in Chapter~4. The purpose of the implementation is not to reproduce a full industrial geophysical simulator, but to provide a research prototype in which several predictive model families can be compared under one physically motivated workflow. In this sense, the system connects user-defined thermoelastic scenarios, effective geological material parameters, model routing, and normalized prediction summaries into one reproducible experimental stack. The implementation therefore acts as the practical bridge between the geological motivation of Chapter~2, the thermoelastic variables and governing relations of Chapter~3, and the comparative prediction task formulated in Chapter~4.

At the architectural level, the implementation is organized as a service-oriented prototype. The full local stack is composed of a backend orchestrator, a web-based scientific frontend, a set of independent model services, a structured catalog of geological media, and a container-based local orchestration layer. In practical terms, the system can be understood as a pipeline in which the user configures a geological medium and a thermoelastic scenario in the frontend, the backend validates and enriches the request, the selected model service performs the prediction, and the normalized response is returned for display, comparison, and summary reporting. This design is appropriate for the thesis because the main scientific question is comparative directional prediction rather than monolithic solver construction.

The frontend is the user-facing layer of the prototype. Its role is to provide a compact scientific interface for selecting a geological medium, choosing a prediction route, specifying the source and probe geometry, defining thermal and mechanical conditions, and visualizing the returned prediction summary. The frontend dynamically loads media and model information from the backend and presents a two-dimensional domain view for interactive experiments. From the perspective of the thesis, the frontend should therefore be described as an experimental configuration and visualization surface rather than as a source of physical modelling. It does not perform the thermoelastic computation itself. Instead, it exposes the structured inputs needed by the prediction workflow and displays the resulting directional quantities and field-level summary indicators.

The backend is the central orchestration layer of the system. This distinction is important: the backend is not the machine-learning model itself. It validates a unified prediction request, resolves the geological medium from the catalog, merges scenario inputs with rock presets, routes the request to the selected remote model service, and normalizes the response before returning it to the caller. The unified request structure carries the selected model, a medium identifier, scenario parameters, source coordinates and direction, probe coordinates, and a domain description. The backend also validates geometric consistency, for example by requiring that the out-of-plane size and resolution collapse to a single planar slice for two-dimensional cases. As a result, the backend provides the common contract that allows different model services to be compared under the same physical input description.

The geological medium catalog is a compact repository of effective material presets. It includes media such as sandstone, limestone, basalt, and granite together with representative properties such as density, total and effective porosity, \(V_p\), \(V_s\), thermal conductivity, heat capacity, and thermal expansion. These entries are not full geological databases and should not be interpreted as field-validated constitutive models. Instead, they serve the thesis purpose of supplying effective material parameters for locally homogeneous and isotropic thermoelastic modelling. This is fully consistent with the simplifications introduced in Chapters~2 and~3, where real geological media are recognized as heterogeneous, porous, fractured, layered, and often anisotropic, but the mathematical formulation proceeds through effective continuum parameters.

The predictive components are separated into distinct services. This separation is useful because it makes the comparison framework explicit and prevents the backend from being mistaken for a single model implementation. The PINN component is the closest to the governing equations of Chapter~3. It predicts the primary fields \((T, u, v)\) and is trained with a combination of a supervised loss, a velocity consistency term, an elastic wave residual, and a coupled thermal residual. The residual terms are built from the thermoelastic stress tensor, displacement gradients, and the coupled heat equation with the \(\gamma T_0 \, \partial \varepsilon_{kk}/\partial t\) term. For this reason, the PINN is the only explicitly physics-informed baseline in the current comparison.

The FNO component is implemented as a dedicated Fourier Neural Operator route and is configured as a trainable two-dimensional baseline. Its current purpose is to learn the evolution of physically meaningful fields on a regular grid, with targets such as temperature and displacement components. Its training objective is a supervised loss of the form
\begin{equation}
\mathcal{L}_{\mathrm{FNO}}
=
\operatorname{MSE}(\widehat{Y},Y)
+ 0.01 \, \operatorname{relativeL2}(\widehat{Y},Y),
\end{equation}
without explicit thermoelastic PDE residual enforcement. It is therefore described in the thesis as a neural-operator baseline rather than as a physics-constrained solver. This distinction matters because the model still operates on physically meaningful fields, yet it does not incorporate the governing thermoelastic equations in the same way as the PINN.

The MeshGraphNet component provides the graph-based route. It can run a MeshGraphNet-style rollout when dataset and checkpoint artifacts are configured, and it can also return a deterministic fallback response when these artifacts are unavailable and fallback is enabled. This means that the component is best described as a graph-based predictive route within a prototype architecture. It is relevant to the thesis because directional propagation in geological media is naturally linked to geometry, local neighborhoods, and mesh-like spatial relationships. At the same time, its practical interpretation must remain careful: depending on the run configuration, the returned result may reflect either a learned graph-based prediction or fallback logic intended to keep the system operational for demonstration and interface testing.

The Transformer-style component is an autoregressive neural-operator baseline parallel to the PINN. Its architecture uses tokens carrying coordinates, current-state fields, and material properties, and predicts \((T, u, v)\) at target coordinates. The current version uses a supervised data loss only and does not yet include physics-informed residuals. It is therefore described in the thesis as an experimental supervised sequence/operator-style baseline whose maturity is lower than that of the main orchestration pattern.

Taken together, these components form a common comparative architecture rather than four physically equivalent solvers. All of them receive inputs motivated by the same thermoelastic problem definition, but they use the physics differently. The PINN is the only component whose training objective explicitly contains thermoelastic residuals derived from the governing equations. The FNO, MeshGraphNet, and Transformer-style routes use physically meaningful inputs and outputs, but they do not enforce the full coupled thermoelastic PDE system in the same way. This wording is important for consistency with Chapter~4 and for avoiding overstatement in the thesis.

Both the mathematical formulation in Chapter~3 and the prototype implementation are two-dimensional. The current experiments operate on planar source--probe paths and two-dimensional cross-sections that are practical for the available data and prototype model components. In such settings, source and probe positions are selected within a single plane, and the directional response is described entirely by the in-plane geometry introduced in Chapter~3.

The deployment structure of the project reinforces this interpretation. The container-based local environment starts the frontend, the backend, and each model component as a separate service with its own health or readiness endpoint. This design is valuable for a bachelor-level research prototype because it isolates component responsibilities, makes failure modes more visible, and allows individual components to be replaced or upgraded without rewriting the entire stack. The project also provides component-specific scripts and procedures for data preparation, training, checkpoint management, validation, and local startup. These support reproducibility and experimentation, even though they do not by themselves establish scientific validation.

In summary, Chapter~6 implements the methodology of Chapter~4 through a backend-centered orchestration stack that connects user inputs, effective geological material presets, model routing, and normalized directional prediction outputs. The implementation should be described as a comparative research prototype with physics-informed and data-driven components, not as a final validated two-dimensional thermoelastic simulator. Its main strengths are unified routing, explicit material handling, component modularity, and transparent model comparison. Its main limitations are equally important: the prototype is two-dimensional, medium presets remain starter values, different model routes use physics at different depths, and fallback-capable components must be interpreted carefully when reporting comparative results.
